\subsection{Literature Revisited}

The work by \cite{gruttmann2000theory} uses \(C^0\) interpolation on the Lagrange basis for \(\alpha\), but employs director interpolation for the section orientation.

\cite{hsiao1999consistent,hsiao2000corotational}

The work by \cite{zhang2011formulation} which derives from \cite{alemdar2001distributed} uses a \(C^1\) field for \(\alpha\). The nonlinear Wagner term is included. The work is limited to symmetric sections.

\cite{rinchen2016formulation}, \cite{rinchen2020geometric} extends the work of \cite{zhang2011formulation} implements the torsion T2/shear T3 formulation in \OpenSees{}. 

\begin{enumerate}
\item \textbf{Torsion Warping}: In this case only a single warping mode is considered so that \(\bm{\alpha}\) is a scalar and the warping mode \(\varphi\) is St. Venant's warping function.
\iffalse
$$
\nabla \omega = \bm{\Phi}\begin{pmatrix}\alpha^{\prime} \\ \alpha \end{pmatrix}
\qquad\text{ with }\qquad
\bm{\Phi} \triangleq \begin{bmatrix}
\varphi & 0 \\ 0 & \varphi_{, y} \\ 0 & \varphi_{, z}
\end{bmatrix}
$$
\fi
In the absence of constraints on \(\alpha\), this corresponds to the RBV formulation described by \cite{armero2021modeling}.
Examples of finite element implementations include \cite{simo1991geometricallyexact} and \cite{gruttmann2000theory}.

$$
  \mathbf{L}_{\mathrm{n}} 
= \mathbf{L}_{\mathrm{m}} 
= \mathbf{L}_{\mathrm{w}} 
= \bm{0}
$$
$$
\TotalStrainMatrix\mathbf{L} = \TotalStrainMatrix = \left[\begin{array}{ccc|ccc|ccc|ccc}
1 & 0 & 0 & 0 & z & - y & \varphi & \psi_{y} & \psi_{z} & 0 & 0 & 0 \\
0 & 1 & 0 & - z & 0 & 0 & 0 & 0 & 0 & \varphi_{,y} & \psi_{y,y} & \psi_{z,y}\\
0 & 0 & 1 & y & 0 & 0 & 0 & 0 & 0 & \varphi_{,z} & \psi_{y,z} & \psi_{z,z}
\end{array}\right]
$$
See Equation~(26) of \cite{dire2018mixed}

% \item For Type 1 torsion and uniform (Type 3) shear warping \eqref{eq:ks-isotropic} reduces to:
% \begin{equation*}
% \hat{\mathbf{K}} = \left[\begin{array}{ccc|ccc|cc}
% E A & 0 & 0 & 0 & E Q_{y} & E Q_{z} & E R_{w} & 0\\
% 0 & G A_{y} & 0 & G Q_{yx} & 0 & 0 & 0 & G R_{y}\\
% 0 & 0 & G A_{z} & G Q_{zx} & 0 & 0 & 0 & G R_{z}\\ \hline
% 0 & G Q_{yx} & G Q_{zx} & G I_{0} & 0 & 0 & 0 & G S_{a} \\
% E Q_{y} & 0 & 0 & 0 & E I_{y} & - E I_{yz} & E S_{y} & 0 \\
% E Q_{z} & 0 & 0 & 0 & - E I_{yz} & E I_{z} & E S_{z} & 0 \\ \hline
% E R_{w} & 0 & 0 & 0 & E S_{y} & E S_{z} & E C_{w} & 0 \\
% 0 & G R_{y} & G R_{z} & G S_{a} & 0 & 0 & 0 & G C_{a}
% \end{array}\right]
% \end{equation*}

\item T2 (constrained) torsion. Referred to as the TWKV formulation by \cite{armero2021modeling}, this formulation imposes the constraint 
  \[
  \alpha = \CurveStrain \cdot \FrameBasis
  \]
  % This assumption derives from observing that, for very thin elements, uniform warping displacements can be considered across the thickness of the cross-section membratures and, thus, transverse shear strains in this direction can be neglected. 
  This assumption is usually adequate for open cross-sections, but can lead to incorrect solutions for closed profiles.
  Examples include \cite{nukala2004mixed}, the \emph{bw3d} element by \cite{battini2002corotational-thesis} and \cite{zhang2011formulation} for symmetric sections, whose work is implemented in \OpenSees{} and derives from \cite{alemdar2001distributed}.

\item \textbf{Linear} kinematics, T2 torsion, and T1 shear

$$
\mathbf{L}_{\mathrm{n}} = \bm{0},
\qquad
\mathbf{L}_{\mathrm{m}} = \left[\begin{matrix}
1 & 0 & 0\\
0 & 0 & 0\\
0 & 0 & 0
\end{matrix}\right],
\quad\text{ and }\quad
\mathbf{L}_{\mathrm{w}} = \left[\begin{matrix}
1 & 0 & 0\\
0 & 0 & 0\\
0 & 0 & 0
\end{matrix}\right]
$$
$$
\TotalStrainMatrix\mathbf{L} = \left[\begin{array}{ccc|ccc|ccc|cc}
1 & 0 & 0 & 0 & z & - y & \varphi & \psi_{y} & \psi_{z} & 0 & 0\\
0 & 1 & 0 & \varphi_{,y} - z & 0 & 0 & 0 & 0 & 0 & \psi_{y,y} & \psi_{z,y}\\
0 & 0 & 1 & \varphi_{,z} + y & 0 & 0 & 0 & 0 & 0 & \psi_{y,z} & \psi_{z,z}\end{array}\right]
$$

\item T2 (constrained) torsion and T2 shear

$$
\TotalStrainMatrix\mathbf{L} = \left[\begin{matrix}
1 & 0 & 0 & \tau r^2 & z & - y & \varphi & \psi_{y} & \psi_{z}\\
0 & \psi_{y,y} + 1 & \psi_{z,y} & \varphi_{,y} - z & 0 & 0 & 0 & 0 & 0\\
0 & \psi_{y,z} & \psi_{z,z} + 1 & \varphi_{,z} + y & 0 & 0 & 0 & 0 & 0\end{matrix}\right]
$$
This adds shear warping to the expression by \cite{battini2002plastic}.

\item T2 (constrained) torsion and T2 shear, linearized

$$
\mathbf{L}_{\mathrm{n}} = \FrameBasis_{\beta} \otimes \FrameBasis_{\beta}, \quad
\mathbf{L}_{\mathrm{m}} = \FrameBasis_{} \otimes \FrameBasis_{},
\quad\text{ and }\quad
\mathbf{L}_{\mathrm{w}} = \bm{1}
$$
% $$
% \mathbf{L}_{\mathrm{p}} = \left[\begin{matrix}
% 0 & 0 & 0 & 1 & 0 & 0\\
% 0 & 1 & 0 & 0 & 0 & 0\\
% 0 & 0 & 1 & 0 & 0 & 0
% \end{matrix}\right]
% \qquad\text{ and }\qquad
% \mathbf{L}_{\mathrm{w}} = \bm{1}
% $$
$$
\TotalStrainMatrix\mathbf{L} = \left[\begin{array}{ccc|ccc|ccc}
1 & 0 & 0 & 0 & z & - y & \varphi & \psi_{y} & \psi_{z}\\
0 & \psi_{y,y} + 1 & \psi_{z,y} & \varphi_{,y} - z & 0 & 0 & 0 & 0 & 0\\
0 & \psi_{y,z} & \psi_{z,z} + 1 & \varphi_{,z} + y & 0 & 0 & 0 & 0 & 0
\end{array}\right]
$$
This is quivalent to Equation~(29) for the ``EVDE'' formulation of  \cite{addessi2021enriched}. 
\item T2 (constrained) torsion and T3 (uniform) shear, linearized

$$
\mathbf{L}_{\mathrm{n}} = \FrameBasis_{\beta} \otimes \FrameBasis_{\beta}, \quad
\mathbf{L}_{\mathrm{m}} = \FrameBasis_{} \otimes \FrameBasis_{},
\quad\text{ and }\quad
\mathbf{L}_{\mathrm{w}} = \bm{1}
$$

$$
\TotalStrainMatrix\mathbf{L} = \left[\begin{array}{ccc|ccc|c}
1 & 0 & 0 & 0 & z & - y & \varphi \\
0 & \psi_{y,y} + 1 & \psi_{z,y} & \varphi_{,y} - z & 0 & 0 & 0 \\
0 & \psi_{y,z} & \psi_{z,z} + 1 & \varphi_{,z} + y & 0 & 0 & 0 
\end{array}\right]
$$
Described by \cite[Section 3.1.2]{lecorvec2012nonlinear} with \(k_y \beta  \triangleq 1+\psi_{y,y}\). This model was also employed by \cite[Equation 3.51]{dire20173d}. 
%
Note that this generalizes the work of \cite{saritas2006mixed}, as presented in Equation~3.23 of \cite{dire20173d}.

\item T3 {Uniform torsion}.  In this case warping is assumed to depend on the torsional component of the curvature $\CurveStrain $ so that $\alpha = \FrameBasis\cdot\CurveStrain $ and $\alpha^{\prime}=0$. Therefore:
$$
\bm{\omega} = (\varphi \FrameBasis\otimes\FrameBasis)\CurveStrain  = \tau \varphi \, \FrameBasis
$$
\begin{Details}
and
$$
\bm{\varepsilon} = \ShearStrain  + \CurveStrain \times\PlanePosition  + \tau\varphi \CurveStrain \times\FrameBasis
$$
\end{Details}

\item Type 3 (uniform) torsion and Type 3 shear warping
$$
\mathbf{L}_{\mathrm{n}} = \FrameBasis_{\beta} \otimes \FrameBasis_{\beta}, \quad
\mathbf{L}_{\mathrm{m}} = \FrameBasis_{} \otimes \FrameBasis_{},
\quad\text{ and }\quad
\mathbf{L}_{\mathrm{w}} = \bm{1}
$$
$$
\TotalStrainMatrix\mathbf{L} = \left[\begin{array}{ccc|ccc}
1 & 0 & 0 & 0 & z & - y \\
0 & \psi_{y,y} + 1 & \psi_{z,y} & \varphi_{,y} - z & 0 & 0 \\
0 & \psi_{y,z} & \psi_{z,z} + 1 & \varphi_{,z} + y & 0 & 0
\end{array}\right]
$$
%
See also \cite{wackerfuss2009mixed}.

\item Type 3 torsion and Type 4 shear  \citep{gruttmann1998geometrical,battini2002corotational-thesis} where additionally \(\psi_{y,y}  \approx \psi_{z,z} \approx 0\)
$$
\TotalStrainMatrix\mathbf{L} = \left[\begin{matrix}
1 & 0 & 0 & 0 & z & - y \\
0 & 1 & 0 & \varphi_{,y} - z & 0 & 0 \\
0 & 0 & 1 & \varphi_{,z} + y & 0 & 0
\end{matrix}\right]
$$

\item Type 4 torsion and Type 3 shear. Described by \cite{lecorvec2012nonlinear} as the \emph{Shear beam} model.
$$
\TotalStrainMatrix\mathbf{L} = \left[\begin{array}{ccc|ccc}
1 & 0 & 0 & 0 & z & - y \\
0 & 1+\psi_{y,y} & 0 & - z & 0 & 0 \\
0 & 0 & 1+\psi_{z,z} &   y & 0 & 0
\end{array}\right]
$$
\item \textbf{Plane Section}: In the absence of warping one has $\bm{\omega}=\bm{0}$ and $\bm{\varepsilon}\equiv \bar{\bm{\varepsilon}}$:
$$
\StrainTensor = \bar{\bm{\varepsilon}}
\stackrel{\mathrm{s}}{\otimes}
\FrameBasis 
+ \frac{1}{2}\left.
\bar{\bm{\varepsilon}}\cdot\bar{\bm{\varepsilon}}\,\FrameBasis\otimes\FrameBasis\right. \\
$$

\end{enumerate}
